%!TEX root = ../main.tex

\begin{abstract}
We develop a theory of endogenous premium formation for Bitcoin treasury
vehicles---publicly traded firms that hold BTC as their primary asset and
issue securities to fund ongoing accumulation.  The NAV premium equals the
present value of a \emph{perpetual accumulation option}: the right to issue
equity at a premium and purchase BTC, thereby increasing BTC-per-share for
existing holders.  The fair premium satisfies a fixed-point equation with a
closed-form steady-state approximation that links it to BTC volatility,
issuance intensity, leverage, and discount rates.  Reflexive feedback between
the premium and accumulation creates multiple equilibria---a high-premium
``flywheel,'' an unstable tipping point, and a low-premium ``death
spiral''---connected by saddle-node bifurcations exhibiting hysteresis.  We
reinterpret five structural indicators (ILE, TEE, PMRI, IBGR, IFRD) as
option Greeks of the embedded accumulation claim: Delta, Adjusted Delta,
Theta, Moneyness, and Exercise Cost.  We derive optimal capital structure
rules governing the transition from ATM equity to convertible debt to
preferred stock as the premium declines, and introduce the Weighted Average
Cost of BTC Acquisition (WACBA).  Market impact analysis yields a reflexivity
amplification factor and conditions for contagion cascades under corporate BTC
adoption.  Calibrating the model to Strategy (ticker: MSTR)---which holds
761,068~BTC (\$56.25B) as of March~2026---we estimate an OU premium process
with mean-reversion speed $\hat\kappa = 5.21$, long-run mean
$\hat\theta = 1.08$, and BTC volatility $\hat\sigma_S = 48.6\%$.  The
current premium of 24.2\% sits below its long-run mean
($\text{PMRI} = -0.71$), implying positive carry and an underpriced
accumulation option.  Three-year survival probability is 91.9\%, and the
implied funding requirement averages \$35.3B.  The framework provides a
unified lens for valuation, risk assessment, and vehicle selection (MSTR vs.\
BTC~ETF).

\medskip
\noindent\textbf{Keywords:} Bitcoin, NAV premium, accumulation option,
reflexivity, multiple equilibria, capital structure, Strategy (MSTR).

\noindent\textbf{JEL Codes:} G12, G13, G32, G23.
\end{abstract}
